\section{The additional geometry for quartics}
\label{sec:quartic-geometry}

In this section $d=4$ and $k=2$.  We retain the number field
$K=\Q(\theta)$, its splitting field $L$, and the distinct roots
$\theta_1,\ldots,\theta_4$.  We first improve the Hilbert bound for
surfaces after the infinite-fibre removal of
Lemma~\ref{lem:quartic-fibers}.  We then construct one fixed exceptional
set of integers that will remove the equality case in the curve argument.

\subsection{The surface bound}

\begin{proposition}\label{prop:quartic-surface}
Let $Y$ be a geometrically irreducible surface defined over $\Q$ in the
rational coordinates of the arithmetic boxes.  Suppose that it meets the open
$a_1a_2a_3a_4\ne0$, satisfies
\[
 a_i^2b_i=n-\theta_i\qquad(1\le i\le4),
\]
and that $n$ is nonconstant and algebraic over the field generated by
$a_1,\ldots,a_4$.
Give the two coordinate blocks positive rational weights $U,V$,
respectively, and put
\[
 D_0=\min(3U,2V,U+V).
\]
Then
\begin{equation}\label{eq:quartic-surface-hilbert}
 H_Y(T)\ge \frac{T^2}{2}
    \min\left\{\frac{(2U+V)D_0}{U^2V^2},\frac3{U^2}\right\}-O(T).
\end{equation}
For fixed weights and a bound on $\deg Y$, the error is uniform in the
moving coordinate changes and defining equations.
\end{proposition}

\begin{proof}
Let $S$ be the closure of the full $a$-image of $Y$.  It is a surface,
because $n$ is algebraic over its function field and the $b_i$ are
$(n-\theta_i)/a_i^2$ on the indicated open.  It is defined over $\Q$ in
the rational coordinates of the first block.

We separate the cases in which the $a$-image has degree at least three,
two, or one.  The first supplies three independent families directly.
A quadric supplies a pair-image bound or an extra field degree; on a
plane, deficient bounds for every pair lead to a contradiction.

We will repeatedly use Proposition~\ref{prop:weighted-hilbert} for
algebraically independent coordinate functions $a_i,a_j$.
Its polynomial $C_1$ has weight $D$ and
involves $b_i$ or $b_j$, since $a_i,a_j$ are independent.  If $D<D_0$,
no monomial can contain two $b$ factors, an $a$ factor and a $b$ factor,
or three $a$ factors.  Thus its equation has the form
\begin{equation}\label{eq:quartic-small-relation}
 c_i b_i+c_j b_j=S_{ij}(a_i,a_j),\qquad
 \deg S_{ij}\le2,\qquad(c_i,c_j)\ne(0,0).
\end{equation}
In particular $D\ge V$.  Substitution of the equations for $b_i,b_j$
gives
\begin{equation}\label{eq:quartic-rational-n}
 \begin{split}
 B_{ij}n&=c_i\theta_i a_j^2+c_j\theta_j a_i^2
                 +a_i^2a_j^2S_{ij}(a_i,a_j),\\
 B_{ij}&=c_i a_j^2+c_j a_i^2.
 \end{split}
\end{equation}
The independence of $a_i,a_j$ implies $B_{ij}\ne0$.  Consequently, in
this case $n,b_i,b_j$ all belong to $\C(a_i,a_j)$.

\paragraph{An $a$-image of degree at least three.}
Suppose $\deg S=e\ge3$.  A general linear projection to an affine plane,
with coordinates $s,t$ linear in the $a_i$, has function-field degree
$e$.  Indeed choose its projective centre at infinity disjoint from the
projective closure of $S$.  A general fibre is a complementary linear
section of degree $e$, with no intersection at infinity; in characteristic
zero its general points are distinct.  A further general linear form $r$
in the $a_i$ separates this finite extension, so $1,r,r^2$ are independent
over $\C(s,t)$.  Products with monomials in $s,t$ give
\[
 H_Y(T)\ge
  \sum_{j=0}^2\#\{(p,q)\in\Z_{\ge0}^2:U(p+q+j)\le T\}
  =\frac{3T^2}{2U^2}-O(T).
\]
This proves the second bound in \eqref{eq:quartic-surface-hilbert}.

\paragraph{A quadric $a$-image.}
Suppose $\deg S=2$.  The degree of an irreducible projective variety is
at least its codimension in its linear span plus one
\cite[Introduction, equation~$(*)$]{EisenbudGreenHulekPopescu2006}.
Applying this to
the projective closure shows that $S$ spans an affine three-space $A$
and is a quadratic hypersurface there: a surface spanning only an affine
plane would equal that plane and have degree one.

The span $A$ is rational in the first-block rational coordinates.  Its
direction is a rational hyperplane in $K$.  The nondegeneracy of the trace
pairing writes its equation as $\operatorname{Tr}_{K/\Q}(\gamma\alpha)=0$
for some $\gamma\in K\setminus\{0\}$.  Write $\ell_i$ for the linear
part of $a_i$ on $A$.  The equation of this direction space is therefore
\[
 \sum_{i=1}^4\sigma_i(\gamma)\ell_i=0,
\]
with all four coefficients nonzero.  In the projective plane of this
direction space, the four lines $\ell_i=0$ have no triple intersection.
Their six pair intersections are therefore distinct.

Let $Q_2$ be the homogeneous quadratic part of the equation of $S$ in
$A$.  It cannot vanish at all six intersections.  Otherwise its
restriction to each of the four lines would have three distinct zeros
and hence vanish identically, forcing all four line equations to divide
the nonzero quadratic $Q_2$.  Choose a pair $i,j$ at whose intersection
direction $v$ one has $Q_2(v)\ne0$.

The fibres of the projection $A\to\mathbb A^2$ given by $(a_i,a_j)$ are
lines parallel to $v$.  On each fibre the equation of $S$ has a nonzero
quadratic leading coefficient $Q_2(v)$.  Thus $S$ has degree two over
this coordinate plane, and $a_i,a_j$ are independent.  If $D\ge D_0$,
Proposition~\ref{prop:weighted-hilbert} already gives the first bound
in \eqref{eq:quartic-surface-hilbert}.  If $D<D_0$,
\eqref{eq:quartic-rational-n} shows that the projected-pair field is
$\C(a_i,a_j)$.  The full field $\C(Y)$ has degree at least two over it,
so the field-extension clause of Proposition~\ref{prop:weighted-hilbert}
doubles the leading coefficient.  Since $D\ge V$ and
$2D\ge2V\ge D_0$, the doubled coefficient suffices as well.

\paragraph{A plane $a$-image.}
It remains to consider $\deg S=1$, so $S=A$ is an affine plane rational
in the first-block coordinates.  Every $a_i$ is nonconstant on $A$:
a nonzero rational direction represents a nonzero element of $K$, and
all its embeddings are nonzero.  Their linear parts span the dual of
the two-dimensional direction space.

Join indices $i,j$ when those linear parts are independent.  The resulting
graph is the complete multipartite graph on their parallelism classes,
with at least two classes.  It is connected and has no isolated vertex.
An edge is exactly an independent pair of affine coordinates on $A$.

Assume for a contradiction that every such pair has $D<D_0$.  Then
\eqref{eq:quartic-small-relation} holds for every edge, and
\eqref{eq:quartic-rational-n} makes $n$ a rational function on $A$.
We distinguish whether it is polynomial.

If $n$ is polynomial, \eqref{eq:quartic-rational-n} shows that its degree
is at most four.  Indeed $B_{ij}$ has degree exactly two for an
independent pair, whereas the right side has degree at most six.
Every index $i$ has nonzero coefficient in at least one relation
\eqref{eq:quartic-small-relation}.  To see this, start with one nonzero
coefficient and apply Galois automorphisms: they preserve $Y$, permute
the conjugate coordinates and roots transitively, and preserve both
independence and the degree bound.  This uses the collection of all
such relations, without requiring any chosen representatives to be
Galois-invariant.

Fix a relation with $c_i\ne0$.  At the generic point of $a_i=0$,
the independent coordinate $a_j$ is a unit.  Polynomiality of $n$
therefore makes $b_j=(n-\theta_j)/a_j^2$ regular there, and the relation
makes $b_i$ regular too.  It follows that
\[
 a_i^2\mid n-\theta_i\quad\hbox{in }\C[A]
 \qquad(1\le i\le4).
\]
The four zero lines are distinct, since their prescribed values
$\theta_i$ are distinct.  Restriction to a general affine line gives
a polynomial of degree at most four with four distinct critical points
and four distinct critical values.  Its derivative would have four
roots despite degree at most three; if the derivative were zero, the
values would all be equal.  This is impossible.

Now suppose $n$ is nonpolynomial.  Both coefficients in every relation
\eqref{eq:quartic-small-relation} must be nonzero, since otherwise it
would express $n$ as a polynomial.  Write $n=P/Q$ in reduced form and
fix an irreducible factor of the nonconstant denominator $Q$.  By
\eqref{eq:quartic-rational-n} it divides every $B_{ij}$.  Over $\C$,
$B_{ij}$ is a product of two distinct affine linear forms
$a_j\pm\zeta a_i$, with $\zeta\ne0$.  The fixed factor therefore defines
a line $\ell$ through the common zero of $a_i,a_j$, transverse to both
of their zero lines.  For each edge, the two intersections with $\ell$
coincide.  Connectivity of the graph forces all four zero lines to
pass through one point $O$.

Center the plane coordinates at $O$, making each $a_i$ homogeneous
linear.  Scale these coordinates by a variable $t$ in
\eqref{eq:quartic-rational-n} and take the rational-function limit at
$t=0$.  For every independent pair the result is the same function
on the projective line of directions:
\begin{equation}\label{eq:quartic-direction-map}
 \mathcal R=\frac{c_i\theta_i a_j^2+c_j\theta_j a_i^2}
          {c_i a_j^2+c_j a_i^2}.
\end{equation}
The term involving $S_{ij}$ acquires an extra factor $t^2$.
The limit is taken in the function field of directions.

For an independent pair, \eqref{eq:quartic-direction-map} is a
nonconstant fractional linear transform of $(a_i/a_j)^2$: its
determinant is nonzero because
$c_i c_j(\theta_j-\theta_i)\ne0$.  It has degree two and is ramified
exactly at the two directions $a_i=0,a_j=0$, with respective values
$\theta_i,\theta_j$.  Since every index occurs in an edge, the same
function would be ramified in four directions with four distinct
values.  Coincident directions already contradict the distinct values,
and four distinct directions contradict its displayed degree-two form.
This proves the required contradiction.

There is consequently an independent pair with $D\ge D_0$, and
Proposition~\ref{prop:weighted-hilbert} proves the first bound in
\eqref{eq:quartic-surface-hilbert} for the plane case.  Together the
three degree cases prove the proposition.  All extra independent
functions used above have bounded weighted degree, so their shifts
preserve the asserted uniform $O(T)$ error.
\end{proof}

\subsection{One global set of quintic values}

The curve argument will encounter degree-five polynomials with four
prescribed critical values.  Their coefficients can initially depend
on the box and on the curve.  The following normalization places all
of them in one fixed finite family before any counting is done.

\begin{proposition}\label{prop:quintic-discard}
There are finitely many polynomials $\Phi\in L[t]$ of degree five with
distinct critical points $z_1=0,z_2=1,z_3,z_4$ satisfying
\[
 \Phi(z_i)=\theta_i\qquad(1\le i\le4).
\]
Let $\mathcal F_f$ be this finite family and put
\[
 \mathcal E_f=
  \{n\in\Z:n=\Phi(z)\text{ for some }
                    \Phi\in\mathcal F_f, z\in L\}.
\]
Then, for $X\ge1$,
\begin{equation}\label{eq:quintic-discard-count}
 \#\bigl(\mathcal E_f\cap[-2X,2X]\bigr)\ll_f X^{1/5}.
\end{equation}
Moreover, if $N(s)\in L[s]$ has degree five and distinct critical
points $\tau_1,\ldots,\tau_4\in L$ with
$N(\tau_i)=\theta_i$, every integer value $N(s)$ at $s\in L$ belongs
to $\mathcal E_f$.
\end{proposition}

\begin{proof}
We first prove finiteness over $\C$.  Introduce the six coefficients
of $\Phi$ and the two variables $z_3,z_4$, impose
$\Phi(z_i)=\theta_i$ and $\Phi'(z_i)=0$, and require that the leading
coefficient and all differences $z_i-z_j$ be nonzero.  This defines a
finite-type parameter locus.  Its four critical points are simple,
because they are distinct and exhaust the roots of the degree-four
polynomial $\Phi'$.

Suppose a reduced irreducible component of this locus had positive
dimension.  Its function field $F/\C$ admits a nonzero derivation
$\partial$ over $\C$: differentiate a member of a transcendence basis
and extend through the finite separable extension.  Write
$Q=\partial\Phi$ for coefficient differentiation.  Differentiating the
value equations gives $Q(z_i)=0$, since $\Phi'(z_i)=0$.  Consequently
\[
 Q=A\Phi',\qquad\deg A\le1.
\]
Differentiating the critical-point equations now gives
\[
 0=Q'(z_i)+\Phi''(z_i)\partial z_i
   =\Phi''(z_i)\bigl(A(z_i)+\partial z_i\bigr).
\]
The second derivative is nonzero at each $z_i$.  The fixed
normalizations $z_1=0,z_2=1$ imply $A(0)=A(1)=0$, so $A=0$.
Thus $\partial$ kills all coefficients of $\Phi$ and all $z_i$.
These generate $F$, contradicting $\partial\ne0$.  The parameter locus
is therefore zero-dimensional and has finitely many complex points.
Its subset of polynomials over $L$ is finite as well.

Fix one such $\Phi$.  Choose an integer $c>0$ such that
\[
 c\Phi(t)=A_5t^5+A_4t^4+\cdots+A_0,
 \qquad A_i\in\mathcal O_L,\quad A_5\ne0.
\]
If $\Phi(z)=n\in\Z$ and $w=A_5z$, then
\begin{equation}\label{eq:quintic-integral-input}
 \begin{split}
 0={}&w^5+A_4w^4+A_3A_5w^3+A_2A_5^2w^2\\
     &+A_1A_5^3w+(A_0-cn)A_5^4.
 \end{split}
\end{equation}
This is monic over $\mathcal O_L$, so $w\in\mathcal O_L$.  Hence every
possible $z$ lies in the fixed fractional lattice
$A_5^{-1}\mathcal O_L$.

For each embedding $\sigma:L\to\C$ one has
$\sigma(\Phi)(\sigma(z))=n$.  The leading term of each of these fixed
degree-five polynomials dominates its lower terms outside a fixed
radius.  For $|n|\le2X$ it follows that
\[
 \max_\sigma|\sigma(z)|\ll_{\Phi,L}X^{1/5}.
\]
Choose a fixed $\Z$-basis $e_1,\ldots,e_{r_L}$ of the fractional lattice,
where $r_L=[L:\Q]$, and write $z=\sum_j u_j e_j$ with $u_j\in\Z$.
Inverting its embedding matrix gives $|u_j|\ll_{\Phi,L}X^{1/5}$ for
every $j$.

We count these lattice points on the locus where $\Phi(z)$ is rational.
To define this locus over $\Q$, expand $\Phi(\sum_j u_je_j)$ in a
rational basis of $L$ beginning with $1$ and set its other coordinates
equal to zero.  After base change to $\C$, the embedding coordinates
identify it with the inverse image of the diagonal under the map
\[
 (w_\sigma)_\sigma\longmapsto
             \bigl(\sigma(\Phi)(w_\sigma)\bigr)_\sigma.
\]
This is a product of finite polynomial maps.  More explicitly, the
inverse image of the diagonal has coordinate ring
\[
 \C[v,(w_\sigma)_\sigma]\big/
       \bigl(\sigma(\Phi)(w_\sigma)-v:\sigma:L\to\C\bigr).
\]
After dividing each relation by its leading coefficient, the equations
are monic in separate variables.  The ring is finite free of rank
$5^{r_L}$ over $\C[v]$, with basis the products of powers
$w_\sigma^{e_\sigma}$ for $0\le e_\sigma<5$.
Thus the rational-value locus has dimension one and degree bounded in
terms of $\Phi,L$.

A bounded-degree affine set of dimension at most one has $O(R)$ integer
points in a coordinate box of radius $R$.  Indeed, on each curve
component choose a nonconstant coordinate, slice at its $O(R)$ integer
values, and apply B\'ezout to each finite fibre.  Zero-dimensional
components contribute only a bounded number of points.  Apply this
bound in the $u_j$ coordinates with $R\ll_{\Phi,L}X^{1/5}$.  It counts
$O_{\Phi,L}(X^{1/5})$ possible $z$, and therefore at most that many
integer values.  Summing over the fixed finite family proves
\eqref{eq:quintic-discard-count}.

Finally, for the polynomial $N$ in the statement set
\[
 \Phi(t)=N\bigl(\tau_1+(\tau_2-\tau_1)t\bigr).
\]
Its coefficients lie in $L$, and its critical points
$(\tau_i-\tau_1)/(\tau_2-\tau_1)$ have the required normalizations
and values.  Hence $\Phi\in\mathcal F_f$.  For every $s\in L$, the
normalized input $(s-\tau_1)/(\tau_2-\tau_1)$ also lies in $L$ and has
the same value under $\Phi$.  This proves the final assertion without
any restriction on the heights of $N$ or its critical points.
\end{proof}
